\subsection{Census-independent completion of the dimension-twenty-four quotient}
\label{sec:direct-d24-quotient}
Let $\mathcal C$ have one of the four based rings
$R9\_1\_\{12,62,87,108\}$, and suppose a central $C_2$ de-equivariantization
has dimension profile $(1^4,2^5)$. Its based ring is necessarily
$R9\_1\_117$. This assertion uses the four source tensors, not a completeness
claim about any catalogue of dimension-twenty-four rings.

Write $z$ for the nontrivial central sign in $\mathcal C$. There is exactly
one possible underlying invertible $z$ with this quotient profile in each
source tensor; the indices are respectively $2,7,3,6$.
The sign orbits determine the restriction matrix $F$ up to relabeling the
quotient simples. Use the quotient basis
\[
 (1,a,b,c,m,x,x',y,y'),\qquad
 d=(1,1,1,1,2,2,2,2,2),\qquad
 \alpha=(b\ c)(x\ x')(y\ y').
\]
The source invertible orbit containing $1$ restricts to $1$, and its other
invertible orbit restricts to $a$. The unique sign-fixed two-dimensional
source object restricts to $b+c$, while the free two-dimensional source orbit
restricts to $m$. The two sign-fixed four-dimensional source objects restrict
to $x+x'$ and $y+y'$. These statements follow from the standard simple-object
description of $C_2$ equivariantization; a fixed quotient simple has two
one-dimensional stabilizer representations, and a free quotient orbit has
one induced simple. Thus this single canonical $F$ covers every actual
restriction functor after changing quotient labels.

Let $C_{ij}^h$ be the source tensor and $n_{ab}^k$ an arbitrary nonnegative
integer quotient tensor. The verifier first enumerates every involution
$s$ on the quotient basis preserving dimension and satisfying
\[
 s(1)=1,\quad s\alpha=\alpha s,\qquad
 F_{i,a}=F_{i^*,s(a)}.
\]
There are eight such involutions for sources $12,62$, and four for
sources $87,108$. For each, introduce all $9^3$ tensor variables and identify
variables only under the necessary equalities
\[
 n_{ij}^k=n_{s(i),k}^{j}=n_{k,s(j)}^{i},\qquad
 n_{ij}^k=n_{\alpha(i),\alpha(j)}^{\alpha(k)}.
\]
The first equalities are Frobenius reciprocity and the last is equivariance.
The resulting number of variables is respectively $92$ or $83$.
Impose the unit and duality equations, and all equations
\[
 \sum_k d_k n_{ij}^k=d_i d_j,
 \qquad
 \sum_{a,b}F_{i,a}F_{j,b}n_{ab}^k
       =\sum_h C_{ij}^h F_{h,k}.
\]
Every equation has nonnegative integer coefficients. For each variable $v$,
an equation containing it with coefficient $w>0$ and right-hand side $r$
gives $v\le\lfloor r/w\rfloor$. The verifier finds and records such a bound
at most one for every variable, before starting the binary search.
In the eight-branch case, sixty variables have bound zero and thirty-two
have bound one; in the four-branch case, the numbers are fifty-two and
thirty-one. Thus the binary search is exhaustive among \emph{all}
nonnegative integer quotient tensors: multiplicity-freeness is an output
of the linear equations, not an assumed census bound.

The finite search is ordinary exhaustive branching with sound propagation.
For an equation with residual target $r$ and unassigned positive weights
$w_1,\ldots,w_t$, the possible sums are computed exactly from
$\prod_j(1+u^{w_j})$ by a bit-set recurrence. A value of a variable is
discarded only when the other weights cannot attain the required residual
sum. If propagation leaves variables unassigned, both values $0,1$ are
examined. Complete linear solutions are checked against every one of the
$9^4$ associativity equations.

For each of sources $12,62$, the complete search has $74$ nodes over all
eight duality branches, producing $36$ linear solutions and exactly four
associative tensors. For each of sources $87,108$, it has eight nodes over
four branches, producing four linear solutions, all associative.
Each resulting tensor is isomorphic to $R9\_1\_117$ under an explicit
unit-preserving permutation. Only at this final identification step is the
$R9\_1\_117$ tensor read; no quotient list or catalogue is consulted during
completion. The receipt records every associative tensor and its
quotient-to-input permutation, every duality branch, source tensor hashes,
all upper bounds and their equation witnesses, and exact search counts.

The standard-library replay is
\nolinkurl{work/resolve_rank9_neargroup/direct_quotient.py}; its full receipt is
\nolinkurl{work/resolve_rank9_neargroup/direct_quotient.json}.
The replay uses explicit error checks and has also been run with Python
optimization enabled.
